\documentclass[../main]{subfiles}
\ifSubfilesClassLoaded{\setcounter{chapter}{3}}{}
\begin{document}

\chapter{Dasar Pemrograman C\# untuk Unity}

\begin{subcpmk}
  \item Sub-CPMK 2.1: Membuat skrip MonoBehaviour dan menggunakan lifecycle functions (Start, Update, Awake, FixedUpdate)
  \item Sub-CPMK 2.2: Menggunakan tipe data Unity (Vector3, Quaternion) dan menerapkan OOP, coroutine, event
\end{subcpmk}

\section{Game Objects}

Setelah memahami jendela Unity Editor (Bab III), bab ini membahas entitas dasar: GameObject dan Component. \unity{GameObject} adalah entitas dasar dalam Unity \cite{unityManual}. Setiap objek di Hierarchy adalah GameObject. GameObject kosong hanya memiliki \class{Transform}; objek 3D (Cube, Sphere, Plane) memiliki Transform plus Mesh Filter dan Mesh Renderer. Mesh Filter menyimpan data geometri; Mesh Renderer menggambar geometri ke layar.

\subsection{Transform}

Transform mengatur posisi, rotasi, dan skala dalam ruang 3D. Koordinat menggunakan sistem tangan kanan: X kanan, Y atas, Z masuk ke layar.

\begin{table}[htbp]
\centering
\begin{tabular}{|l|p{6cm}|}
\hline
\textbf{Properti} & \textbf{Deskripsi} \\
\hline
Position & Koordinat X, Y, Z dalam ruang world. Origin (0,0,0) di pusat scene \\
\hline
Rotation & Rotasi dalam derajat (Euler). Urutan: X, Y, Z \\
\hline
Scale & Ukuran relatif. (1,1,1) = normal. (2,2,2) = dua kali lipat \\
\hline
\end{tabular}
\caption{Properti Transform}
\end{table}

\subsection{Parent-Child}

GameObject dapat menjadi child dari GameObject lain. Transform child relatif terhadap parent. Memindahkan parent akan memindahkan semua child. Berguna untuk mengelompokkan objek (misal: semua collectible dalam satu parent).

\subsection{Primitive 3D}

Unity menyediakan primitive: Cube, Sphere, Cylinder, Plane, Capsule, Quad. Untuk Roll a Ball, kita gunakan Plane (lantai), Sphere (player), dan Cube (tembok arena).

\section{Variabel, Tipe Data, dan Kontrol Alur}

\begin{learningoutcomes}
\item Menggunakan tipe data dasar C\# dan tipe Unity (Vector3, Quaternion)
\item Menggunakan public, private, dan SerializeField
\item Menerapkan if-else, switch, loop dalam konteks game
\end{learningoutcomes}

\subsection{Tipe Data Dasar dan Unity}

\textbf{Dasar}: int, float, bool, string, char. \textbf{Unity}: Vector3, Vector2, Quaternion, Color, GameObject, Transform.

\begin{lstlisting}[language={[Sharp]C}]
using UnityEngine;

public class VariableExample : MonoBehaviour
{
    public int health = 100;
    public float speed = 5.5f;
    public Vector3 position;
    
    [SerializeField] private int maxHealth = 100;
}
\end{lstlisting}

\subsection{Kontrol Alur}

\begin{lstlisting}[language={[Sharp]C}]
// If-else
if (health > 0) { /* Player hidup */ }
else { GameOver(); }

// Switch
switch (gameState) {
    case GameState.Menu: ShowMenu(); break;
    case GameState.Playing: StartGame(); break;
}

// Loop
for (int i = 0; i < 10; i++) Debug.Log(i);
foreach (GameObject enemy in enemies) enemy.SetActive(false);
\end{lstlisting}

\begin{catatan}
Public variables muncul di Inspector. Gunakan SerializeField untuk private variables yang perlu diatur di Inspector.
\end{catatan}

\section{Object-Oriented Programming dalam Unity}

\begin{learningoutcomes}
\item Memahami penggunaan class dan inheritance di Unity
\item Menggunakan virtual/override untuk polymorphism
\item Menerapkan encapsulation
\end{learningoutcomes}

\subsection{Class dan Inheritance}

\begin{lstlisting}[language={[Sharp]C}]
public class BaseEnemy : MonoBehaviour
{
    protected int health;
    
    public virtual void Attack()
    {
        Debug.Log("Enemy attacks");
    }
}

public class MeleeEnemy : BaseEnemy
{
    public override void Attack()
    {
        Debug.Log("Melee attack");
    }
}
\end{lstlisting}

\subsection{Encapsulation}

Gunakan private untuk data internal, public methods untuk aksi. Gunakan protected untuk anggota yang bisa diakses subclass.

\section{Coroutine dan Event}

\begin{learningoutcomes}
\item Memahami konsep coroutine dan IEnumerator
\item Menggunakan yield return untuk delay
\item Mengimplementasikan event system
\end{learningoutcomes}

\subsection{Coroutine}

\begin{definisi}[Coroutine]
Coroutine adalah fungsi yang dapat menjeda eksekusinya sendiri dan melanjutkan di frame berikutnya \cite{UnityScriptingAPI2024}.
\end{definisi}

\begin{lstlisting}[language={[Sharp]C}]
using System.Collections;
using UnityEngine;

IEnumerator Countdown()
{
    for (int i = 5; i > 0; i--)
    {
        Debug.Log(i);
        yield return new WaitForSeconds(1f);
    }
}
\end{lstlisting}

\subsection{Event dan Delegate}

\begin{lstlisting}[language={[Sharp]C}]
using System;

public static event Action OnGameOver;
public static event Action<int> OnScoreChanged;

// Subscribe di OnEnable, Unsubscribe di OnDisable
void OnEnable() { GameManager.OnGameOver += HandleGameOver; }
void OnDisable() { GameManager.OnGameOver -= HandleGameOver; }
\end{lstlisting}

\begin{catatan}
Selalu unsubscribe event di OnDisable untuk menghindari memory leaks.
\end{catatan}


\begin{aktivitas}
  \item \textbf{Skrip Pertama}: Buat skrip C\# yang menampilkan pesan di Start() dan Update(). Lampirkan ke GameObject.
  \item \textbf{Lifecycle}: Buat skrip yang menampilkan urutan Awake, Start, Update, OnDestroy dengan Debug.Log.
  \item \textbf{Coroutine}: Buat countdown 3-2-1-Go menggunakan coroutine.
  \item \textbf{Event}: Implementasikan event OnScoreChanged dan subscribe dari script lain.
\end{aktivitas}

\begin{latihan}
  \item Jelaskan perbedaan Awake(), Start(), dan Update()!
  \item Kapan menggunakan FixedUpdate() vs Update()?
  \item Apa fungsi SerializeField? Kapan menggunakannya?
  \item Jelaskan cara kerja coroutine dengan yield return!
  \item Mengapa harus unsubscribe event di OnDisable?
\end{latihan}

\begin{asesmen}
\textbf{Instrumen Penilaian untuk Sub-CPMK 2.1, 2.2}

\textbf{A. Pilihan Ganda}: Mengenai MonoBehaviour, lifecycle, tipe data Unity, coroutine.
\textbf{B. Praktik}: Buat skrip Player dengan health, TakeDamage(), dan event OnDeath.
\textbf{Rubrik}: Lihat Lampiran A
\end{asesmen}

\begin{checklist}
  \item Saya dapat membuat skrip MonoBehaviour
  \item Saya memahami lifecycle functions
  \item Saya dapat menggunakan Vector3, Transform, dll
  \item Saya dapat menggunakan coroutine
  \item Saya dapat mengimplementasikan event
\end{checklist}

\begin{rangkuman}
Bab ini membahas dasar C\# untuk Unity: MonoBehaviour, lifecycle (Awake, Start, Update, FixedUpdate), tipe data (Vector3, Quaternion), kontrol alur, OOP, coroutine, dan event system.
\end{rangkuman}

\ifSubfilesClassLoaded{
  \renewcommand{\bibname}{Daftar Pustaka}
  \bibliographystyle{plain}
  \bibliography{../references}
}{}
\end{document}
